\section{Exact localization and finite corrections with both diffusion coefficients}
\label{sec:localized-layer}

We now reconstruct the localization identities and finite correction mechanism
of Alp\"oge and Buckmaster, \emph{Blowup for the Boussinesq equations with
smooth forcing}, Sections 4 and 6, keeping the coordinates and frequency
parameters of Section~\ref{sec:boussinesq-core}. We also derive the
corresponding finite correction equations when the thermal diffusion
coefficient $\kappa$ and the viscosity $\nu$ are nonnegative. This is an
exact finite-layer calculation. The limiting many-layer construction and
uniform force bounds up to a singular time are not asserted here.

\subsection{Material coordinates and the full two residuals}

Let $[t_0,t_1]$ be a compact time interval. Retain the smooth trace-free
matrix $D(t)$ and vector $G(t)$, the backgrounds $u_b=Dx$ and
$\theta_b=G\cdot x$, and their prescribed forces from
\eqref{eq:bc-base-scalar}--\eqref{eq:bc-base-vector}.
Solve
\[
M'=DM,\qquad M(t_0)=\mathrm{Id},\qquad
p=\zeta(t_0)\ne0,\qquad \zeta=M^{-T}p.
\]
The existence and invertibility proof in Section~\ref{sec:boussinesq-core}
applies with initial time $t_0$. Differentiating the determinant by its
polynomial cofactor formula gives
$(\det M)'=(\operatorname{tr}D)\det M=0$, so $\det M=1$.
Define, without changing units,
\begin{gather*}
a=M^{-1}x,\qquad s=\lambda\zeta\cdot x,\qquad C=M^{-1}M^{-T},\\
q=p\cdot Cp=|\zeta|^2,\qquad
\widetilde G=M^TG,\qquad c=M^{-1}e_1.
\end{gather*}
The phase variable $s$ and material variable $a$ are independent variables
of a profile, even though its evaluation on physical space ties them
together. For a smooth $2\pi$-periodic profile $A(s,a,t)$ put
\[
A^\natural(x,t)=A(\lambda\zeta(t)\cdot x,M(t)^{-1}x,t).
\]
All profile time derivatives below fix $(s,a)$. Introduce the operators
\begin{align*}
\partial_\perp&=Jp\cdot\nabla_a,&
B&=(Cp)\cdot\nabla_a,&
\Delta_C&=\operatorname{div}_a(C\nabla_a),\\
D_1&=2B\partial_s,&
L_\lambda&=\lambda^2q\partial_s^2+\lambda D_1+\Delta_C,&
\mathcal P&=\lambda\zeta\partial_s+M^{-T}\nabla_a,\\
J_s(A,B_0)&=A_s\partial_\perp B_0-(B_0)_s\partial_\perp A,&
\{A,B_0\}_a&=J\nabla_aA\cdot\nabla_aB_0.
\end{align*}
Here $B_0$ in the last line is an arbitrary profile, not the operator $B$.
The original source denotes $q$ by $\kappa$; we retain $q$ so that
$\kappa$ can continue to mean thermal diffusivity.

The exact chain rules are
\begin{align}
\nabla_xA^\natural&=(\mathcal PA)^\natural,&
(\partial_t+Dx\cdot\nabla_x)A^\natural&=(A_t)^\natural,
\label{eq:ll-chain-time}\\
\Delta_xA^\natural&=(L_\lambda A)^\natural,&
\nabla_x^\perp\Psi^\natural\cdot\nabla_xA^\natural
&=\bigl(\lambda J_s(\Psi,A)+\{\Psi,A\}_a\bigr)^\natural.
\label{eq:ll-chain-space}
\end{align}
For the first line, differentiate $a=M^{-1}x$ and
$s=\lambda\zeta\cdot x$ and use $\zeta'=-D^T\zeta$.
Squaring the spatial gradient gives the three terms in $L_\lambda$.
For the last identity, expand both gradients. In dimension two a
matrix of determinant one satisfies $M^{-1}JM^{-T}=J$, as direct
multiplication by the four entries shows. Thus the slow-slow term is
$\{\Psi,A\}_a$, the fast-fast term is zero, and the two cross terms
are respectively $\lambda\Psi_s\partial_\perp A$ and
$-\lambda A_s\partial_\perp\Psi$. This proves every sign in
\eqref{eq:ll-chain-space}.

For arbitrary smooth profiles $V,\Psi$, add
\[
\theta=\theta_b+V^\natural,\qquad
u=u_b+v,\qquad v=\nabla_x^\perp\Psi^\natural,
\qquad \omega=\omega_b+(L_\lambda\Psi)^\natural,
\quad \omega_b=D_{21}-D_{12}.
\]
The exact increments to scalar force and curl of vector force are
\begin{align}
E_\theta={}&V_t+\lambda(J\zeta\cdot G)\Psi_s
 +J\nabla_a\Psi\cdot\widetilde G
 +\lambda J_s(\Psi,V)+\{\Psi,V\}_a-\kappa L_\lambda V,
\label{eq:ll-full-scalar}\\
E_\omega={}&\partial_t(L_\lambda\Psi)
 +\lambda J_s(\Psi,L_\lambda\Psi)+\{\Psi,L_\lambda\Psi\}_a
 -\lambda\zeta_1V_s-c\cdot\nabla_aV-\nu L_\lambda^2\Psi.
\label{eq:ll-full-vorticity}
\end{align}
Indeed $v\cdot\nabla\theta_b=v\cdot G$ gives the two displayed
background-gradient terms in \eqref{eq:ll-full-scalar} because
$JM^{-T}=MJ$. The background vorticity is spatially constant, so
$v\cdot\nabla\omega_b=0$. Applying
\eqref{eq:ll-chain-time}--\eqref{eq:ll-chain-space} to every other
term of the scalar and vorticity equations gives the formulas exactly.
In particular, $\partial_t(L_\lambda\Psi)$ differentiates $q$ and $C$.

These curl residuals have a direct full momentum realization. Put
$\mathsf v=J\mathcal P\Psi$ and define the profile vector
\begin{equation}
\mathsf E_u=2D\mathsf v+J\mathcal P\Psi_t
 +(\mathsf v\cdot\mathcal P)\mathsf v
 -\nu J\mathcal P L_\lambda\Psi-Ve_2.
\label{eq:ll-vector-force}
\end{equation}
Then with the unchanged pressure $p_b$ the full momentum equation
holds for $f_u=f_{u,b}+\mathsf E_u^\natural$, and
$\operatorname{curl}_x\mathsf E_u^\natural=E_\omega^\natural$.
To verify the vector formula, differentiating the coefficients of
$J\mathcal P$ gives $D J\mathcal P$: for example
$J\zeta'=DJ\zeta$, and $J(M^{-T})'=DJM^{-T}$.
It follows that
\[
(\partial_t+Dx\cdot\nabla_x)v
=Dv+(J\mathcal P\Psi_t)^\natural.
\]
The additional term $(v\cdot\nabla)u_b=Dv$ accounts for the second
copy of $Dv$. The self-advection, diffusion, and buoyancy increments
are exactly the remaining terms of \eqref{eq:ll-vector-force}.
For the curl assertion, expansion in coordinates gives, for any smooth
divergence-free two-dimensional $u$,
\[
\operatorname{curl}((u\cdot\nabla)u)
=u\cdot\nabla(\partial_1u_2-\partial_2u_1)
 +(\partial_1u_1+\partial_2u_2)
  (\partial_1u_2-\partial_2u_1)
=u\cdot\nabla\omega.
\]
Taking the curl of the already proved vector equation therefore gives
\eqref{eq:ll-full-vorticity}. For compact material support, both force
increments and the added fields have compact spatial support on this
finite interval. The prescribed affine background separately retains
the spatial behavior stated in Section~\ref{sec:boussinesq-core}.

\subsection{The complete leading localized residual}

Retain the odd periodic $F$, its even mean-zero primitive $P$, and
$P'=F$. Allow arbitrary smooth signed amplitudes $T(t),Q(t)$ and an
arbitrary smooth material envelope $g(a,t)$. Set
\[
V_0=TFg,\qquad \Psi_0=QP g,\qquad
\mathcal O=\lambda^2qQ,\qquad h=Bg,\qquad d=\Delta_Cg.
\]
Products here and below are evaluated at $(s,a,t)$ only at the end.
Direct application of $L_\lambda$ gives
\begin{align}
L_\lambda V_0&=T\bigl(\lambda^2qF''g+2\lambda F'h+Fd\bigr),
\label{eq:ll-lap-scalar}\\
L_\lambda\Psi_0&=\mathcal O F'g+2\lambda QFh+QP d,
\label{eq:ll-lap-stream}\\
L_\lambda^2\Psi_0&=Q\bigl[
 \lambda^4q^2F'''g+4\lambda^3qF''h
 +\lambda^2F'(2qd+4Bh)+4\lambda F Bd+P\Delta_Cd\bigr].
\label{eq:ll-bilap-stream}
\end{align}
The last line follows by applying the three terms of $L_\lambda$
to each of the three terms in \eqref{eq:ll-lap-stream}.
Its two mixed second-order contributions are
$2\lambda^3qF''h$ each, its two $F'd$ contributions are
$\lambda^2qF'd$ each, and
$\Delta_Ch=B\Delta_Cg=Bd$ because $C$ and $Cp$ are constant
in $a$. These observations give exactly the coefficients $4,2,4,4$
in \eqref{eq:ll-bilap-stream}, with no term omitted.

Substitution in \eqref{eq:ll-full-scalar} gives
\begin{align}
E_\theta(V_0,\Psi_0)={}&
 (T'+\lambda(J\zeta\cdot G)Q)Fg+TFg_t
 +QP(J\nabla_ag\cdot\widetilde G)
\notag\\
&+\lambda QT(F^2-PF')g\partial_\perp g
 -\kappa T\bigl(\lambda^2qF''g+2\lambda F'h+Fd\bigr).
\label{eq:ll-expanded-scalar}
\end{align}
The complete vorticity formula is
\begin{align}
E_\omega(V_0,\Psi_0)={}&
 (\mathcal O'-\lambda\zeta_1T)F'g+\mathcal O F'g_t
 +\partial_t(2\lambda QFh+QP d)-TF(c\cdot\nabla_ag)
\notag\\
&+\lambda^3qQ^2(FF'-PF'')g\partial_\perp g
\notag\\
&+2\lambda^2Q^2\bigl(F^2g\partial_\perp h-PF'h\partial_\perp g\bigr)
 +2\lambda Q^2PF\{g,h\}_a
\notag\\
&+\lambda Q^2PF\bigl(g\partial_\perp d-d\partial_\perp g\bigr)
 +Q^2P^2\{g,d\}_a
\notag\\
&-\nu Q\bigl[
 \lambda^4q^2F'''g+4\lambda^3qF''h
 +\lambda^2F'(2qd+4Bh)+4\lambda F Bd+P\Delta_Cd\bigr].
\label{eq:ll-expanded-vorticity}
\end{align}
For a verification of all nonlinear terms, the slow-slow bracket of
$QP g$ with $TFg$ vanishes, and their fast bracket is
$QT(F^2-PF')g\partial_\perp g$.
In the vorticity equation, bracket $QP g$ successively with the three
summands of \eqref{eq:ll-lap-stream}. The first gives
$\lambda^3qQ^2(FF'-PF'')g\partial_\perp g$ and zero slow bracket.
The second gives
\[
2\lambda^2Q^2(F^2g\partial_\perp h-PF'h\partial_\perp g)
\quad\hbox{and}\quad 2\lambda Q^2PF\{g,h\}_a.
\]
The third gives
$\lambda Q^2PF(g\partial_\perp d-d\partial_\perp g)$ and
$Q^2P^2\{g,d\}_a$. This exhausts the two brackets and proves
\eqref{eq:ll-expanded-vorticity}.
Time differentiation in that formula keeps the derivatives of $C$,
$q$, $g$, $Q$, $h$, and $d$.

For the exact inviscid amplitudes of \eqref{eq:bc-amplitudes}, a smooth
activation $w(t)$, and a time-independent material envelope, put
\[
T=w\Theta,\qquad Q=\frac{w\Omega}{\lambda^2q}.
\]
Then the first terms in \eqref{eq:ll-expanded-scalar} and
\eqref{eq:ll-expanded-vorticity} are respectively $w'\Theta Fg$ and
$w'\Omega F'g$. Setting $\kappa=\nu=0$ therefore reproduces every
term of the source's equation (6.5). The compact support is generated
by the material envelope; no spatial coordinate or frequency has been
set equal to one in this calculation.

Define phase average at fixed $(a,t)$ by
\[
\langle A\rangle(a,t)=\frac1{2\pi}\int_0^{2\pi}A(s,a,t)\,ds,
\qquad \mu_2=\langle F^2\rangle,\qquad \mu_P=\langle P^2\rangle.
\]
The exact means, including arbitrary $g_t$, $\kappa$, and $\nu$, are
\begin{equation}
\langle E_\theta\rangle=2\mu_2\lambda QTg\partial_\perp g,
\qquad
\langle E_\omega\rangle
=2\mu_2\lambda^2Q^2\partial_\perp(gh)
 +\mu_PQ^2\{g,d\}_a.
\label{eq:ll-phase-means}
\end{equation}
Indeed $\langle F\rangle=\langle P\rangle=0$ and all derivatives
have zero mean. Integration by parts gives
$\langle PF'\rangle=-\mu_2$,
$\langle FF'\rangle=\langle PF''\rangle=0$, and
$\langle PF\rangle=0$. Substitution of these identities in the two
expanded residuals proves \eqref{eq:ll-phase-means} term by term.
These averages are functions of material position. They are not
physical spatial integrals and need not vanish.

\subsection{A finite phase correction that includes viscosity}

We prove an exact finite cancellation identity. It extends the source's
inviscid phase recursion; the construction here treats diffusion in the
evolution equation itself. It does not replace viscosity by an
uncontrolled forcing term.

Fix $\lambda>0$, $\kappa,\nu\geq0$ and the smooth coefficients above
on $[t_0,t_1]$. Since $q>0$, compactness gives a positive lower bound
for $q$ on this interval. Let $g\in C_c^\infty(\mathbb R^2)$ be
time-independent and equal to one on an open material set $U$.
For the parity assertion below, additionally let $g(-a)=g(a)$.
Let $w$ be smooth and flat at $t_0$, meaning $w^{(j)}(t_0)=0$ for
every nonnegative integer $j$.
Set
\[
K=\lambda^2q,\qquad b=J\zeta\cdot G,\qquad
\alpha=\frac{b}{\lambda q},\qquad \beta=\lambda\zeta_1.
\]
Here $K$ is an actual coefficient, not a change of spatial or temporal
units. Start with any smooth mean-zero odd initial phase pair and solve
\begin{equation}
\partial_t\binom{\mathcal T}{\mathcal W}
=\begin{pmatrix}\kappa K\partial_s^2&-\alpha\\
                   \beta&\nu K\partial_s^2\end{pmatrix}
  \binom{\mathcal T}{\mathcal W}.
\label{eq:ll-phase-evolution}
\end{equation}
The existence and smoothness proof is given below. Define the leading
profiles and streamfunctions by
\begin{equation}
V_0=wg\mathcal T,\qquad W_0=wg\mathcal W,
\qquad
\Psi_n=K^{-1}\partial_s^{-1}W_n\quad(n\geq0).
\label{eq:ll-level-def}
\end{equation}
The inverse $\partial_s^{-1}$ is the unique mean-zero periodic
primitive of a mean-zero profile. Explicitly, for such a profile $H$,
\[
\partial_s^{-1}H(s,a,t)
=\int_0^sH(r,a,t)\,dr
 -\frac1{2\pi}\int_0^{2\pi}\int_0^vH(r,a,t)\,dr\,dv.
\]
The zero average makes the first integral periodic. Differentiation
and averaging verify respectively that its derivative is $H$ and its
mean is zero. Two such primitives differ by a constant with mean zero,
so uniqueness follows. All these operations fix $a$ and preserve its
support and vanishing on $U$.

Write $\mathsf P_0H=H-\langle H\rangle$. All negatively indexed
profiles are zero. For $j\geq0$ define
\begin{equation}
Z_j=\partial_sW_j+\lambda D_1\Psi_{j-1}+\Delta_C\Psi_{j-2},
\qquad Z_j^\flat=Z_j-\partial_sW_j.
\label{eq:ll-z-def}
\end{equation}
For $n\geq1$, the following sources involve only profiles of index
less than $n$:
\begin{align}
S_n^\theta={}&J\nabla_a\Psi_{n-1}\cdot\widetilde G
 +\lambda\sum_{i+j=n-1}J_s(\Psi_i,V_j)
 +\sum_{i+j=n-2}\{\Psi_i,V_j\}_a
\notag\\
&-\kappa\bigl(\lambda D_1V_{n-1}+\Delta_CV_{n-2}\bigr),
\label{eq:ll-source-theta}\\
S_n^\omega={}&\partial_tZ_n^\flat-c\cdot\nabla_aV_{n-1}
 +\lambda\sum_{i+j=n-1}J_s(\Psi_i,Z_j)
 +\sum_{i+j=n-2}\{\Psi_i,Z_j\}_a
\notag\\
&-\nu\bigl(K\partial_s^2Z_n^\flat
       +\lambda D_1Z_{n-1}+\Delta_CZ_{n-2}\bigr).
\label{eq:ll-source-omega}
\end{align}
The sums have nonnegative indices; an empty sum is zero. In particular,
$Z_n^\flat$ contains only $\Psi_{n-1}$ and $\Psi_{n-2}$, so it
does not introduce an unknown of level $n$ into its own source.

For each $n\geq1$ solve the zero-data equation
\begin{equation}
\partial_t\binom{V_n}{W_n}
=\begin{pmatrix}\kappa K\partial_s^2&-\alpha\\
                   \beta&\nu K\partial_s^2\end{pmatrix}
 \binom{V_n}{W_n}
 +\binom{-\mathsf P_0S_n^\theta}
              {-\partial_s^{-1}\mathsf P_0S_n^\omega},
\qquad (V_n,W_n)|_{t=t_0}=0.
\label{eq:ll-correction-evolution}
\end{equation}
This is a fully specified recursion at every finite depth; no
unsolved equation is substituted for its solution. The explicit
solution follows from the two-dimensional propagators described next.

For $m\in\mathbb Z$ let $\Phi_m(t,\tau)$ solve
\[
\partial_t\Phi_m(t,\tau)=A_m(t)\Phi_m(t,\tau),\qquad
\Phi_m(\tau,\tau)=\mathrm{Id},\qquad
A_m(t)=\begin{pmatrix}-\kappa K(t)m^2&-\alpha(t)\\
                           \beta(t)&-\nu K(t)m^2\end{pmatrix}.
\]
An explicit convergent definition for $t\geq\tau$ is the ordered
integral series
\begin{equation}
\Phi_m(t,\tau)=\mathrm{Id}
 +\sum_{r=1}^\infty
 \int_{\tau\leq t_r\leq\cdots\leq t_1\leq t}
 A_m(t_1)\cdots A_m(t_r)\,dt_r\cdots dt_1.
\label{eq:ll-propagator-series}
\end{equation}
For fixed $m$, if $L_m=\sup\|A_m\|$, the $r$th summand has norm
at most $L_m^r(t-\tau)^r/r!$. Uniform convergence and the fundamental
theorem of calculus give the matrix differential equation.
Uniqueness follows by iterating the integral equation for a difference,
as in Section~\ref{sec:boussinesq-core}.

For a profile $H$ write
$\widehat H_m=(2\pi)^{-1}\int_0^{2\pi}H(s)e^{-ims}\,ds$.
Put $\widehat R_{n,0}=0$ and, for $m\ne0$,
\[
\widehat R_{n,m}(a,t)=
\binom{-\widehat S_{n,m}^\theta(a,t)}
      {-(im)^{-1}\widehat S_{n,m}^\omega(a,t)}.
\]
Then the required solution is exactly
\begin{equation}
\binom{V_n}{W_n}(s,a,t)
=\sum_{m\ne0}e^{ims}
 \int_{t_0}^t\Phi_m(t,\tau)\widehat R_{n,m}(a,\tau)\,d\tau.
\label{eq:ll-duhamel}
\end{equation}
The initial pair in \eqref{eq:ll-phase-evolution} has the analogous
series with coefficient $\Phi_m(t,t_0)$ times its initial Fourier
coefficient. For $\kappa=\nu=0$, $A_m$ is independent of $m$, so
\eqref{eq:ll-duhamel} becomes the pointwise phase Duhamel integral
of the same $2\times2$ amplitude system as in the source.

Here are convergence details needed to justify this formula as a
smooth solution, including the degenerate cases where a diffusion
coefficient is zero. For a complex solution $z'=A_mz$, its squared
Euclidean norm satisfies
\[
\frac{d}{dt}|z|^2
=-2\kappa Km^2|z_1|^2-2\nu Km^2|z_2|^2
 +2\operatorname{Re}\bigl((\beta-\alpha)\overline z_2z_1\bigr)
\leq |\beta-\alpha||z|^2.
\]
Multiplication by
$\exp(-\int_\tau^t|\beta-\alpha|)$ and differentiation therefore
give
\begin{equation}
\|\Phi_m(t,\tau)\|\leq
\exp\left(\frac12\int_\tau^t|\beta(r)-\alpha(r)|\,dr\right),
\label{eq:ll-propagator-bound}
\end{equation}
uniformly in $m$. Smooth periodic sources have Fourier coefficients
decaying faster than every inverse power of $|m|$: integration by
parts $N$ times bounds each coefficient by $|m|^{-N}$ times the
supremum of its $N$th phase derivative. The same argument applies
after any fixed material or time derivative, uniformly on the compact
time interval and compact material support. Bound
\eqref{eq:ll-propagator-bound} therefore proves uniform convergence
of \eqref{eq:ll-duhamel} after arbitrarily many phase or material
derivatives. Its first time derivative is obtained by the fundamental
theorem of calculus and the matrix equation; multiplication by
$A_m$ loses at most two powers of $m$. Repetition proves convergence
of every mixed time derivative, since coefficient derivatives are
bounded on the compact interval and every finite power loss is
available from source smoothness. Induction on $n$ now constructs
all levels as smooth profiles. Complex conjugation of the Fourier
coefficients for $m$ and $-m$ proves that the resulting profiles
are real. Their zero phase means follow either from the displayed
sum or from the homogeneous zero-mode equation with zero initial data.
This proves existence and the explicit formula. Fourier expansion
and uniqueness of the matrix ODE at each frequency prove uniqueness.

\begin{proposition}[Exact finite residual with two diffusion coefficients]
\label{prop:ll-finite}
For any integer $N\geq0$, construct levels $0,\ldots,N$ as above
and put
\[
V=\sum_{j=0}^NV_j,\qquad \Psi=\sum_{j=0}^N\Psi_j.
\]
All levels $j\geq1$ vanish on $U$ and outside $\operatorname{supp}g$
for every phase. They have zero phase means and are flat at $t_0$.
If $g$ is even, $V_j$ and $W_j$ are odd, and $\Psi_j$ is even,
under $(s,a)\mapsto(-s,-a)$.

Set the primary profiles $V_j,W_j,\Psi_j$ equal to zero for $j>N$,
while retaining the definition \eqref{eq:ll-z-def}; thus
\[
Z_{N+1}=\lambda D_1\Psi_N+\Delta_C\Psi_{N-1},\qquad
Z_{N+2}=\Delta_C\Psi_N,\qquad Z_j=0\quad(j>N+2).
\]
Use these truncated profiles in the source formulas for all positive
indices. The exact residuals are
\begin{align}
E_\theta&=w'g\mathcal T
 +\sum_{n=1}^N\langle S_n^\theta\rangle
 +\sum_{n=N+1}^{2N+2}S_n^\theta,
\label{eq:ll-finite-scalar}\\
E_\omega&=w'g\partial_s\mathcal W
 +\sum_{n=1}^N\langle S_n^\omega\rangle
 +\sum_{n=N+1}^{2N+4}S_n^\omega.
\label{eq:ll-finite-vorticity}
\end{align}
Every summand is explicitly determined by the finite Duhamel formula;
no infinite-layer convergence statement is included in these identities.
\end{proposition}

\begin{proof}
At level zero the profiles are independent of $a$ on $U$ and vanish
outside $\operatorname{supp}g$. Every term in each positive-level
source contains a material derivative: this is explicit in the linear
terms, in $Z_n^\flat$, and in both brackets. If every earlier positive
level vanishes on $U$, its derivatives vanish there. The only possible
zeroth-level factors are independent of $a$ there, so the entire source
vanishes. The same argument holds outside the support. Phase averaging,
phase inversion, and the propagator integral do not change $a$.
Induction proves both support assertions.

The factor $w$ is flat initially, so the leading levels and all their
mixed derivatives are zero at $t_0$. The source formulas preserve this
property through sums, products, spatial derivatives, and time
derivatives. In \eqref{eq:ll-correction-evolution}, differentiating
successively in time at $t_0$ shows that each higher profile and all
its mixed derivatives are zero there. This proves flatness.

For parity, a derivative in either $s$ or a material coordinate
reverses simultaneous parity. Thus $L_\lambda$ preserves parity,
both brackets have the product parity of their two arguments,
$Z_j$ is even when $W_j$ is odd and $\Psi_j$ is even,
$S_n^\theta$ is odd, and $S_n^\omega$ is even.
The phase average has the same simultaneous parity as its original
profile, and phase inversion reverses it. Consequently both source
components in \eqref{eq:ll-correction-evolution} are odd. The
equation and its unique zero-data solution preserve this parity.
Since $\Psi_j=K^{-1}\partial_s^{-1}W_j$, it is even.
Induction begins with the odd initial phase pair and even $g$.

It remains to prove the residual equalities, retaining the time
dependence of $K$. The definition of $\Psi_n$ gives the identity
$K\partial_s^2\Psi_n=\partial_sW_n$ for every $t$. Therefore
the principal scalar expression at level $n$ is
\[
\partial_tV_n+\alpha W_n-\kappa K\partial_s^2V_n,
\]
and the principal vorticity expression is
\[
\partial_s\bigl(\partial_tW_n-\beta V_n
                         -\nu K\partial_s^2W_n\bigr).
\]
For $n\geq1$, equation \eqref{eq:ll-correction-evolution} makes
these $-\mathsf P_0S_n^\theta$ and
$-\mathsf P_0S_n^\omega$, respectively. At $n=0$,
\eqref{eq:ll-phase-evolution} leaves precisely $w'g\mathcal T$
and $w'g\partial_s\mathcal W$.

To account for all remaining terms, the finite index shift in
\eqref{eq:ll-z-def} gives
\[
L_\lambda\Psi=\sum_{j=0}^{N+2}Z_j.
\]
The two scalar brackets give grades $i+j+1$ and $i+j+2$,
the background-gradient and mixed-diffusion terms give grades one
and two higher than their primary profile, and hence all are
exactly the terms in \eqref{eq:ll-source-theta}.
Their largest possible grade is $2N+2$.
In the vorticity equation, after removing the principal expression,
the linear time derivative is $\partial_tZ_n^\flat$, and the
remaining viscosity terms are exactly
\[
-\nu\bigl(K\partial_s^2Z_n^\flat
                  +\lambda D_1Z_{n-1}+\Delta_CZ_{n-2}\bigr).
\]
The two vorticity brackets give grades $i+j+1$ and $i+j+2$ with
$i\leq N$ and $j\leq N+2$. Thus all remaining terms are exactly
\eqref{eq:ll-source-omega}, with maximal grade $2N+4$.
Adding each source to its canceled mean-zero part leaves its phase
mean through grade $N$ and the whole source at the higher grades.
This is precisely \eqref{eq:ll-finite-scalar}--\eqref{eq:ll-finite-vorticity}.
No differentiation of $q$ has been omitted: it was incorporated in
the exact identity $K\partial_s^2\Psi_n=\partial_sW_n$ before
differentiating.
\end{proof}

The affine core in the source can now be stated exactly. In the
inviscid case, take
$\mathcal T=\Theta(t)F(s)$ and $\mathcal W=\Omega(t)F(s)$,
where \eqref{eq:bc-amplitudes} holds. These solve
\eqref{eq:ll-phase-evolution}. On $a\in U$, every finite corrected
field therefore agrees with the uncut leading field. If also
$F(s)=s$ for $|s|<s_0$, then on
\[
\{x:M^{-1}x\in U,\ |\lambda\zeta\cdot x|<s_0\}
\]
one has, with the original signs and coefficients,
\[
\vartheta=w\Theta\lambda\zeta\cdot x,\qquad
v=\frac{w\Omega}{q}(J\zeta\otimes\zeta)x,\qquad
\varpi=w\Omega.
\]
The activation terms $w'g\mathcal T$ and
$w'g\partial_s\mathcal W$ are still present there; the corrections
to the state vanish, not necessarily the force.
For a sine phase pair and nonzero diffusion, the damped amplitude
system \eqref{eq:bc-diffusive-amplitudes} gives an explicit solution
of \eqref{eq:ll-phase-evolution}, and the same material plateau is
preserved. For a general profile initially linear near zero, the
diffusive phase evolution is given instead by its Fourier propagator;
the preceding proof does not replace that evolution by the inviscid
linear-core identities. In every case the finite residual is the
explicit quantity \eqref{eq:ll-finite-scalar}--\eqref{eq:ll-finite-vorticity},
including its phase means and final uncanceled grades.
